\chapter{Logic Gates and Boolean Algebra}
\label{ch:gates}

In Chapter~\ref{ch:signals} the door sensor became one bit and the LED
another, and between them stood a single inverter. But the security panel needs a
more complex rule: raise the alarm if the door is open \emph{and} (it is night \emph{or}
the guard mode is on). How do we tell this to a circuit? For that we need logic
functions, logic gates, and Boolean algebra, by which expressions can be
simplified and rewritten.

\begin{buildbox}
Open Logra, go to the \emph{Learn} tab, and complete the lessons ``NAND from blocks'' and
``De Morgan's law''. Then open the \code{gates} example and the
\emph{Waveforms} tab: go through the input combinations 00, 01, 10, 11 and find
the row where exclusive OR differs from ordinary OR.
\end{buildbox}

\section{Logic Functions}

A \term{logic function} assigns to each combination of input
bits an output value: 0 or 1. A function is conveniently specified by a \term{truth
table}---a list of all input combinations and the output values. A function of
$n$ inputs has $2^n$ rows in its table.

A circuit that computes a logic function is called a
\term{logic gate}. Logic gates are drawn with
standard symbols. In Russia the GOST standard is used (a rectangle with a
symbol inside), while in foreign literature and in software the
American (ANSI) symbols are often used. It is useful to recognize both.

Functions from everyday life are everywhere. A hallway lamp with two
switches: the lamp is on if the switches are in \emph{different} positions
(this is exclusive OR). An elevator moves only if the doors are closed \emph{and} the
button is pressed (this is AND). A doorbell rings from the button at the building entrance \emph{or} from the button
at the apartment door (this is OR).

\begin{diybox}[The alarm in words and as a table]
\textbf{Task.} The panel rule: raise the alarm if the door is open and (it is night or
the guard mode is on). Write it as a truth table: inputs door, night, armed.\par
\textbf{In Logra.} The \code{logic} example (the \emph{Designer} and
\emph{Waveforms} tabs)---see how several inputs are combined into an
expression.\par
\textbf{Expected result.} A table of eight rows; the output is one in three
rows: door${}=1$ and at least one of night, armed equals 1.\par
\textbf{How to check.} Check the number of ones: $2^3=8$ rows, of which three have
an alarm.
\begin{itemize}[leftmargin=3.4em,labelwidth=3em,labelsep=0.4em,align=left]
  \item[\stars{1}] Repeat: build the table for $F = A\cdot B + \overline{A}\cdot C$
    (problem~1).
  \item[\stars{2}] Modify: add a fourth input (a panic button that always triggers the alarm)
    and count the rows of the new table.
  \item[\stars{3}] Project: describe in words and as a table the rule ``the siren sounds if
    there is an alarm and the user has not entered the code''; how many rows does the table have?
\end{itemize}
\end{diybox}

\section{Basic Gates}

\newcommand{\gaterow}[5]{%
  \begin{tikzpicture}[baseline=-1mm]
    \node[#1] (g) at (0,0) {};
    \ifstrequal{#1}{gNOT}{\draw[wire] (g.input) -- ++(-0.35,0);}{%
      \draw[wire] (g.input 1) -- ++(-0.35,0); \draw[wire] (g.input 2) -- ++(-0.35,0);}
    \draw[wire] (g.output) -- ++(0.35,0);
  \end{tikzpicture} &
  \begin{tikzpicture}[baseline=-1mm]
    \node[#2] (h) at (0,0) {};
    \ifstrequal{#2}{iNOT}{\draw[wire] (h.input) -- ++(-0.35,0);}{%
      \draw[wire] (h.input 1) -- ++(-0.35,0); \draw[wire] (h.input 2) -- ++(-0.35,0);}
    \draw[wire] (h.output) -- ++(0.35,0);
  \end{tikzpicture} & #3 & #4 & #5 \\[6mm]}

\begin{table}[H]
\centering
\caption{Basic logic gates. The ``table'' column gives the output values
for the inputs $ab = 00, 01, 10, 11$}
\label{tab:gates_all}
\small
\begin{tabular}{c c l l c}
\toprule
\textbf{ANSI} & \textbf{GOST / IEC} & \textbf{Name} & \textbf{Formula} & \textbf{Table} \\
\midrule
\gaterow{gNOT}{iNOT}{NOT (inversion)}{$y=\overline{a}$}{\code{10} (for $a=0,1$)}
\gaterow{gAND}{iAND}{AND (conjunction)}{$y=a\cdot b$}{\code{0001}}
\gaterow{gOR}{iOR}{OR (disjunction)}{$y=a+b$}{\code{0111}}
\gaterow{gNAND}{iNAND}{NAND (Sheffer stroke)}{$y=\overline{a\cdot b}$}{\code{1110}}
\gaterow{gNOR}{iNOR}{NOR (Peirce arrow)}{$y=\overline{a+b}$}{\code{1000}}
\gaterow{gXOR}{iXOR}{Exclusive OR (XOR)}{$y=a\oplus b$}{\code{0110}}
\gaterow{gXNOR}{iXNOR}{Exclusive NOR (XNOR)}{$y=\overline{a\oplus b}$}{\code{1001}}
\bottomrule
\end{tabular}
\end{table}

A small circle on the output of a gate denotes inversion. That is why NAND is drawn as
an AND gate with a circle, and NOR as an OR gate with a circle.

In words:
\begin{itemize}
  \item \textbf{NOT}: the output is the opposite of the input;
  \item \textbf{AND}: the output is 1 only if \emph{all} inputs are 1;
  \item \textbf{OR}: the output is 1 if \emph{at least one} input is 1;
  \item \textbf{Exclusive OR}: the output is 1 if the inputs \emph{differ}
        (for many inputs---if the number of ones is odd).
\end{itemize}

\begin{figure}[H]
\centering
\begin{tikztimingtable}[timing/slope=0.05, xscale=2.4, yscale=1.4, timing/font=\sffamily\small, timing/rowdist=3.2ex]
  $a$              & 2L 2L 2H 2H \\
  $b$              & 2L 2H 2L 2H \\
  $a\cdot b$       & 6L 2H \\
  $a+b$            & 2L 6H \\
  $\overline{a\cdot b}$ & 6H 2L \\
  $a\oplus b$      & 2L 4H 2L \\
\end{tikztimingtable}
\caption{The gates on a timing diagram: all four
input combinations are tried in turn}
\end{figure}

\begin{diybox}[Basic gates in action]
\textbf{Task.} Build the NOR and exclusive NOR (equality) gates and
verify that their tables match the table in the chapter.\par
\textbf{In Logra.} The \emph{Learn} tab, the lessons ``NOR from blocks'' and
``Equality: XNOR''; the \code{gates} example on the \emph{Designer} tab.\par
\textbf{Expected result.} NOR gives \code{1000}, XNOR gives \code{1001} for
the inputs 00, 01, 10, 11.\par
\textbf{How to check.} The built-in lesson check goes through all
combinations; on the \emph{Waveforms} tab compare the output with the table.
\begin{itemize}[leftmargin=3.4em,labelwidth=3em,labelsep=0.4em,align=left]
  \item[\stars{1}] Repeat: complete both lessons.
  \item[\stars{2}] Modify: in the \code{gates} example find the combination where exclusive
    OR differs from OR, and explain why.
  \item[\stars{3}] Project: build the function $F=A\cdot B+\overline{A}\cdot C$ from
    problem~1. Acceptance criterion: all eight combinations on the waveform
    match your truth table.
\end{itemize}
\end{diybox}

\section{How a Gate Is Built Inside}

Inside chips, logic gates are built from transistors. In
modern CMOS technology, pairs of transistors of two types are used:
a $p$-channel transistor conducts when its gate is at zero, an $n$-channel one when it is at one.
The simplest gate---the inverter---consists of just two transistors.

\begin{figure}[H]
\centering
\begin{circuitikz}[scale=0.95, transform shape]
  \draw (0,0) node[pmos, anchor=D] (P) {};
  \draw (P.S) node[vcc] {$+3.3$\,V};
  \draw (P.D) -- ++(0,-0.3) coordinate (mid) -- ++(0,-0.3) node[nmos, anchor=D] (N) {};
  \draw (N.S) node[ground] {};
  \draw (P.G) -- ++(-0.5,0) coordinate (g1);
  \draw (N.G) -- ++(-0.5,0) coordinate (g2);
  \draw (g1) -- (g2);
  \draw ($(g1)!0.5!(g2)$) to[short, -o] ++(-1,0) node[left] {$a$};
  \draw (mid) to[short, *-o] ++(1.6,0) node[right] {$y=\overline{a}$};
  \node[font=\sffamily\small, align=left, anchor=west] at (3.6,0.2) {$a=0$: the upper one conducts, $y$ is connected to $+3.3$\,V $\to$ 1};
  \node[font=\sffamily\small, align=left, anchor=west] at (3.6,-1.9) {$a=1$: the lower one conducts, $y$ is connected to ground $\to$ 0};
\end{circuitikz}
\caption{A CMOS inverter: two transistors work as switches}
\end{figure}

NAND and NOR gates are made of four transistors, while AND and OR are
a NAND and a NOR with an inverter on the output (six transistors). That is why in
chips the inverting gates turn out to be the ``cheapest''.

\begin{infobox}[Why 3.3~V]
The first mass-produced series of logic chips (TTL, the 1960s) ran on
5~V, and this voltage remained the standard for a long time. Later transistors became
so small that their very thin gate layer cannot withstand 5~V.
Besides, at every switching the circuit spends energy of about
$C\,U^2$, where $C$ is the load capacitance and $U$ is the supply voltage: by lowering the
voltage from 5 to 3.3~V we get $(3.3/5)^2\approx 0.44$, that is, almost
half the losses. That is how 3.3~V appeared, then 2.5, 1.8 and 1.2~V; the core of
modern chips runs at about 1~V. The thresholds of 0.8 and 2.0~V in 3.3~V logic
were kept the same as in TTL, so that old and new chips could be
connected directly.
\end{infobox}

\begin{diybox}[What a gate is made of]
\textbf{Task.} Build a NOT gate from NAND gates by connecting both inputs
together, and count the transistors of the CMOS versions: NOT, NAND, AND.\par
\textbf{In Logra.} The \emph{Learn} tab, the lesson ``NAND from blocks'' and the lesson
``XOR from four NANDs''; the \emph{Designer} tab.\par
\textbf{Expected result.} NOT takes 2 transistors, NAND 4, AND 6;
exclusive OR built from NANDs requires four gates.\par
\textbf{How to check.} The built-in lesson check; for your own circuit,
go through the inputs on the \emph{Waveforms} tab.
\begin{itemize}[leftmargin=3.4em,labelwidth=3em,labelsep=0.4em,align=left]
  \item[\stars{1}] Repeat: build NOT from NAND (problem~4).
  \item[\stars{2}] Modify: build XOR from four NANDs and compare with your calculation of
    the number of gates.
  \item[\stars{3}] Project: estimate how many times more transistors an
    exclusive OR built from NANDs needs than a single NAND.
\end{itemize}
\end{diybox}

\section{Boolean Algebra}

Logic expressions can be handled almost like ordinary
algebraic ones: expand brackets, factor out common terms, cancel.
These rules are the subject of \term{Boolean algebra}, named after the
English mathematician George Boole. The AND operation is written as multiplication
($a\cdot b$ or $ab$), OR as addition ($a+b$), and NOT with an overbar
($\overline{a}$).

\begin{table}[H]
\centering
\caption{The basic laws of Boolean algebra}
\label{tab:bool}
\begin{tabular}{l l l}
\toprule
\textbf{Law} & \textbf{For AND} & \textbf{For OR} \\
\midrule
Operations with constants & $a\cdot 0 = 0$, \ $a\cdot 1 = a$ & $a+1 = 1$, \ $a+0 = a$ \\
Idempotence            & $a\cdot a = a$ & $a+a = a$ \\
Complement             & $a\cdot\overline{a} = 0$ & $a+\overline{a} = 1$ \\
Double negation        & \multicolumn{2}{c}{$\overline{\overline{a}} = a$} \\
Commutative            & $ab = ba$ & $a+b = b+a$ \\
Associative            & $(ab)c = a(bc)$ & $(a+b)+c = a+(b+c)$ \\
Distributive           & $a(b+c) = ab+ac$ & $a+bc = (a+b)(a+c)$ \\
Absorption             & $a(a+b) = a$ & $a+ab = a$ \\
\textbf{De Morgan's laws} & $\overline{a\cdot b} = \overline{a}+\overline{b}$ & $\overline{a+b} = \overline{a}\cdot\overline{b}$ \\
\bottomrule
\end{tabular}
\end{table}

\begin{warnbox}[Warning: an unusual law]
The second distributive law $a+bc = (a+b)(a+c)$ does not hold in ordinary
arithmetic, but it is true in Boolean algebra. Check it with a truth table!
\end{warnbox}

\subsection{De Morgan's Laws}

De Morgan's laws are the most useful in practice. They say: \emph{an inversion
over an expression can be ``pushed'' inside by replacing AND with OR and vice versa}.
Graphically this means that a NAND gate can be drawn as an OR gate with
inversions on its inputs:

\begin{figure}[H]
\centering
\begin{tikzpicture}
  \node[gNAND] (g1) at (0,0) {};
  \draw[wire] (g1.input 1) -- ++(-0.5,0) node[left] {$a$};
  \draw[wire] (g1.input 2) -- ++(-0.5,0) node[left] {$b$};
  \draw[wire] (g1.output) -- ++(0.6,0) node[right] {$\overline{ab}$};
  \node[font=\Large] at (3.6,0) {$=$};
  \begin{scope}[xshift=7.4cm]
    \node[gOR] (g2) at (0,0) {};
    \node[gNOT] (n1) at (-2.4,0.7) {};
    \node[gNOT] (n2) at (-2.4,-0.7) {};
    \draw[wire] (n1.output) -| ([xshift=-4mm]g2.input 1) -- (g2.input 1);
    \draw[wire] (n2.output) -| ([xshift=-4mm]g2.input 2) -- (g2.input 2);
    \draw[wire] (n1.input) -- ++(-0.6,0) node[left] {$a$};
    \draw[wire] (n2.input) -- ++(-0.6,0) node[left] {$b$};
    \draw[wire] (g2.output) -- ++(0.6,0) node[right] {$\overline{a}+\overline{b}$};
  \end{scope}
\end{tikzpicture}
\caption{De Morgan's law as a circuit}
\end{figure}

\begin{table}[H]
\centering
\caption{Checking De Morgan's law with a truth table}
\begin{tabular}{cc|c|cc|c}
\toprule
$a$ & $b$ & $\overline{a\cdot b}$ & $\overline{a}$ & $\overline{b}$ & $\overline{a}+\overline{b}$ \\
\midrule
0 & 0 & \textbf{1} & 1 & 1 & \textbf{1} \\
0 & 1 & \textbf{1} & 1 & 0 & \textbf{1} \\
1 & 0 & \textbf{1} & 0 & 1 & \textbf{1} \\
1 & 1 & \textbf{0} & 0 & 0 & \textbf{0} \\
\bottomrule
\end{tabular}
\end{table}

\begin{diybox}[Two circuits for one function]
\textbf{Task.} Prove De Morgan's law in practice: build two circuits,
$\overline{a\cdot b}$ and $\overline{a}+\overline{b}$, and compare their outputs.\par
\textbf{In Logra.} The \emph{Learn} tab, the lesson ``De Morgan's law''; the example
\code{demorgan} (the \emph{Designer} and \emph{Waveforms} tabs).\par
\textbf{Expected result.} The outputs of both circuits coincide on all four
combinations; the ``mismatch'' output is always 0.\par
\textbf{How to check.} The \emph{Waveforms} tab: go through the input combinations.
\begin{itemize}[leftmargin=3.4em,labelwidth=3em,labelsep=0.4em,align=left]
  \item[\stars{1}] Repeat: complete the lesson ``De Morgan's law''.
  \item[\stars{2}] Modify: build and compare $\overline{a+b}$ and
    $\overline{a}\cdot\overline{b}$ (problem~3).
  \item[\stars{3}] Project: simplify $A+A\cdot B$ using the absorption law and verify on the
    waveform that the result matches $A$ (problem~2).
\end{itemize}
\end{diybox}

\section{Universal Gates}

A surprising fact: \textbf{any logic circuit can be built from NAND gates
alone}. The same is true for NOR. Such gates are called
\term{universal}. It is enough to show how to get NOT, AND and OR from NAND---and
everything else can be built from them.

\begin{figure}[H]
\centering
\begin{adjustbox}{max width=\linewidth}
\begin{tikzpicture}
  % NOT
  \node[font=\sffamily\small\bfseries] at (0.3,1.3) {NOT};
  \node[gNAND] (a1) at (0.3,0) {};
  \draw[wire] (a1.input 1) -- ++(-0.4,0) coordinate (t1);
  \draw[wire] (a1.input 2) -- ++(-0.4,0) coordinate (t2);
  \draw[wire] (t1) -- (t2); \draw[wire] ($(t1)!0.5!(t2)$) -- ++(-0.6,0) node[left] {$a$};
  \fill ($(t1)!0.5!(t2)$) circle (1.3pt);
  \draw[wire] (a1.output) -- ++(0.5,0) node[right] {$\overline{a}$};
  % AND
  \begin{scope}[xshift=4.3cm]
  \node[font=\sffamily\small\bfseries] at (1.3,1.3) {AND};
  \node[gNAND] (b1) at (0,0.0) {};
  \node[gNAND] (b2) at (2.6,0) {};
  \draw[wire] (b1.input 1) -- ++(-0.5,0) node[left] {$a$};
  \draw[wire] (b1.input 2) -- ++(-0.5,0) node[left] {$b$};
  \draw[wire] (b2.input 1) -- ++(-0.3,0) coordinate (u1);
  \draw[wire] (b2.input 2) -- ++(-0.3,0) coordinate (u2);
  \draw[wire] (u1) -- (u2);
  \draw[wire] (b1.output) -- ($(u1)!0.5!(u2)$);
  \fill ($(u1)!0.5!(u2)$) circle (1.3pt);
  \draw[wire] (b2.output) -- ++(0.5,0) node[right] {$ab$};
  \end{scope}
  % OR
  \begin{scope}[xshift=10.6cm]
  \node[font=\sffamily\small\bfseries] at (1.3,1.6) {OR};
  \node[gNAND] (c1) at (0,0.8) {};
  \node[gNAND] (c2) at (0,-0.8) {};
  \node[gNAND] (c3) at (2.6,0) {};
  \foreach \c/\n in {c1/a, c2/b} {
    \draw[wire] (\c.input 1) -- ++(-0.4,0) coordinate (u1);
    \draw[wire] (\c.input 2) -- ++(-0.4,0) coordinate (u2);
    \draw[wire] (u1) -- (u2); \draw[wire] ($(u1)!0.5!(u2)$) -- ++(-0.5,0) node[left] {$\n$};
    \fill ($(u1)!0.5!(u2)$) circle (1.3pt);
  }
  \draw[wire] (c1.output) -| ([xshift=-4mm]c3.input 1) -- (c3.input 1);
  \draw[wire] (c2.output) -| ([xshift=-4mm]c3.input 2) -- (c3.input 2);
  \draw[wire] (c3.output) -- ++(0.5,0) node[right] {$a+b$};
  \end{scope}
\end{tikzpicture}
\end{adjustbox}
\caption{NOT, AND and OR built only from NAND gates}
\end{figure}

Let us check the OR circuit using De Morgan's laws:
$\overline{\overline{a}\cdot\overline{b}} = \overline{\overline{a}} + \overline{\overline{b}} = a + b$.

\begin{infobox}[Why this matters]
In the era of the K155 series chips (TTL), designers often built entire
devices from K155LA3 chips, each containing four NAND gates.
This made the circuit cheaper and required fewer different parts. \lookahead{in
programmable chips (FPGAs) the question is solved differently: there, instead of gates,
truth tables are used (Chapter~\ref{ch:inside}).}
\end{infobox}

\begin{diybox}[Everything from NAND]
\textbf{Task.} Build the NOT, AND and OR gates using only NAND gates.
Then build the panel rule $\text{led\_n}=\overline{\text{door}\cdot
(\text{night}+\text{armed})}$ entirely from NAND gates.\par
\textbf{In Logra.} The \emph{Learn} tab, the lesson ``NAND from blocks''; the
\emph{Designer} tab.\par
\textbf{Expected result.} NOT takes 1 gate, AND 2, OR 3; the panel
rule comes out with four NAND gates.\par
\textbf{How to check.} The \emph{Waveforms} tab: go through the eight combinations of
door, night, armed; the output \code{led\_n} is 0 (the LED is on) in three
rows.
\begin{itemize}[leftmargin=3.4em,labelwidth=3em,labelsep=0.4em,align=left]
  \item[\stars{1}] Repeat: build AND and OR from NAND.
  \item[\stars{2}] Modify: build the expression without the common inversion, using De Morgan's
    law, and compare the number of gates.
  \item[\stars{3}] Project: build the circuit with four NANDs and confirm the
    table from the first section (problem~10).
\end{itemize}
\end{diybox}

\section{Multi-Input Gates}

The AND, OR and exclusive OR gates can have more than two inputs. They can be
built from two-input gates by connecting them in a ``chain'' or in a ``tree'':

\begin{figure}[H]
\centering
\begin{tikzpicture}
  \node[font=\sffamily\small\bfseries] at (2.6,2.6) {Chain: 3 levels};
  \node[gAND] (a1) at (0,1.7) {};
  \node[gAND] (a2) at (2.6,0.85) {};
  \node[gAND] (a3) at (5.2,0) {};
  \draw[wire] (a1.input 1) -- ++(-0.5,0) node[left] {$a$};
  \draw[wire] (a1.input 2) -- ++(-0.5,0) node[left] {$b$};
  \draw[wire] (a1.output) -| ([xshift=-4mm]a2.input 1) -- (a2.input 1);
  \draw[wire] (a2.input 2) -- (a2.input 2 -| a1.input 1) -- ++(-0.5,0) node[left] {$c$};
  \draw[wire] (a2.output) -| ([xshift=-4mm]a3.input 1) -- (a3.input 1);
  \draw[wire] (a3.input 2) -- (a3.input 2 -| a1.input 1) -- ++(-0.5,0) node[left] {$d$};
  \draw[wire] (a3.output) -- ++(0.5,0) node[right] {$abcd$};
  \begin{scope}[xshift=9.4cm]
  \node[font=\sffamily\small\bfseries] at (1.3,2.6) {Tree: 2 levels};
  \node[gAND] (b1) at (0,1.1) {};
  \node[gAND] (b2) at (0,-1.1) {};
  \node[gAND] (b3) at (2.6,0) {};
  \draw[wire] (b1.input 1) -- ++(-0.5,0) node[left] {$a$};
  \draw[wire] (b1.input 2) -- ++(-0.5,0) node[left] {$b$};
  \draw[wire] (b2.input 1) -- ++(-0.5,0) node[left] {$c$};
  \draw[wire] (b2.input 2) -- ++(-0.5,0) node[left] {$d$};
  \draw[wire] (b1.output) -| ([xshift=-4mm]b3.input 1) -- (b3.input 1);
  \draw[wire] (b2.output) -| ([xshift=-4mm]b3.input 2) -- (b3.input 2);
  \draw[wire] (b3.output) -- ++(0.5,0) node[right] {$abcd$};
  \end{scope}
\end{tikzpicture}
\caption{A four-input AND gate built from two-input gates. The tree is faster:
the signal passes through two gates rather than three}
\end{figure}

\begin{ideabox}
The seven basic logic gates (NOT, AND, OR, NAND, NOR, exclusive
OR and exclusive NOR) are defined by truth tables. Expressions with them are simplified by
the laws of Boolean algebra, and De Morgan's laws let us change AND to OR.
The NAND and NOR gates are universal: any circuit can be built from either of them.
\end{ideabox}

\begin{diybox}[Alarm: stage 2 of the project]
\textbf{Task.} Build the security panel rule
$\text{alarm}=\text{door}\cdot(\text{night}+\text{armed})$ and show it on an
LED.\par
\textbf{In Logra.} The \emph{Learn} tab, the lesson ``AND with three inputs''; the example
\code{exactly-two} as a model of multi-input logic; the
\emph{Designer} tab.\par
\textbf{Expected result.} The alarm is 1 in three combinations out of eight:
$(\text{door},\text{night},\text{armed})=101,\ 110,\ 111$.\par
\textbf{How to check.} The \emph{Waveforms} tab: go through all eight
combinations and compare the output with the table from the first section; on the
\emph{RTL} tab compare the circuit with the figure in the chapter.
\begin{itemize}[leftmargin=3.4em,labelwidth=3em,labelsep=0.4em,align=left]
  \item[\stars{1}] Repeat: complete the lesson ``AND with three inputs''.
  \item[\stars{2}] Modify: in the \code{exactly-two} example explain how ``exactly two''
    differs from ``at least two''.
  \item[\stars{3}] Project. \emph{Stage of the ``Security Panel'' project:
    $\text{alarm}=\text{door}\,\&\,(\text{night}\,|\,\text{armed})$.} Acceptance
    criterion: out of the eight combinations the alarm is on in exactly three, and the board
    LED (which lights on zero) turns on in exactly those.
\end{itemize}
\end{diybox}

\section{Verilog Examples}
\label{sec:gates-verilog}

Logic gates are written in Verilog almost exactly as in Boolean algebra. The
five examples below are complete modules from the folder
\code{code/b2\_gates/}; each has a testbench \code{tb\_*.v} that sweeps all the
input combinations and prints \code{OK}. Running them in Logra: the \emph{Code}
tab $\to$ \emph{Synthesize} (\code{Ctrl+Enter}) $\to$ \emph{RTL} (the circuit
of gates) $\to$ \emph{Waveforms} (an input sweep); see
Section~\ref{sec:signals-verilog} for details. The language itself is covered in
Chapter~\ref{ch:verilog}.

\subsection{Example 1. All the Gates}

Every element of Table~\ref{tab:gates_all} corresponds to one Verilog operator.

\vfile{\codepath/b2_gates/all_gates.v}{code/b2\_gates/all\_gates.v}

\textbf{How to read it.}
\begin{itemize}
  \item One module with two inputs and seven outputs; each \code{assign} line is one gate.
  \item The operators: \code{\textasciitilde} is NOT, \code{\&} is AND, \code{|} is OR, \code{\^{}} is exclusive OR. The inverting variants (NAND, NOR, XNOR) are obtained by inverting the result.
  \item Parentheses fix the order of evaluation: \code{\textasciitilde(a \& b)} is NAND, while \code{\textasciitilde a \& b} is a quite different function (invert first, then AND).
\end{itemize}
\textbf{Checking.} \code{tb\_all\_gates.v} prints the truth table of all the
gates for $ab=00,01,10,11$ and compares it with the table of the chapter
(\code{0001} for AND, \code{0111} for OR, and so on). In Logra look at the
circuit on the \emph{RTL} tab and at the same tables over time on the
\emph{Waveforms} tab. Try to find the row where exclusive OR differs from OR.

\textbf{Finish it yourself.} Add an ``inhibit'' gate: the output is 1 when
$a=1$ and $b=0$.

\textbf{Sample answer.}\par\nopagebreak
\begin{vcode}
module inhibit (
    input  a, b,
    output y          // a AND NOT b
);
    assign y = a & ~b;
endmodule
\end{vcode}

\subsection{Example 2. The Universality of NAND}

From NAND gates alone we build NOT, AND, OR and exclusive OR. In Verilog this is
especially clear: there is one small ``brick'' module, and everything else is
its instances.

\vfile{\codepath/b2_gates/nand_universal.v}{code/b2\_gates/nand\_universal.v}

\textbf{How to read it.}
\begin{itemize}
  \item \code{module nand\_gate} is the brick: \code{y = \textasciitilde(a \& b)}. Everything else is built only from its instances.
  \item \code{nand\_gate g1 (.a(a), .b(a), .y(y));} is an instance of the module with named connections: the module's port \code{.a} is connected to the outside wire \code{a}. Inputs tied together turn the NAND into a NOT.
  \item \code{wire n, p, q;} are internal wires between the gates. The XOR of four NANDs repeats the circuit from the chapter's problem: a shared node \code{n} and three gates around it.
\end{itemize}
\textbf{Checking.} \code{tb\_nand\_universal.v} connects the four circuits to the
same inputs and compares them with the built-in operators \code{\textasciitilde},
\code{\&}, \code{|}, \code{\^{}} on all four combinations. In Logra synthesize
\code{xor\_from\_nand} as the top module and compare its circuit on \emph{RTL}
with the expression \code{a \^{} b}.

\textbf{Finish it yourself.} Build a NOR using only the operation ``NAND''
(\code{\textasciitilde(x \& y)}).

\textbf{Sample answer.}\par\nopagebreak
\begin{vcode}
module nor_from_nand (
    input  a, b,
    output y
);
    wire na = ~(a & a);        // NOT a
    wire nb = ~(b & b);        // NOT b
    wire o  = ~(na & nb);      // OR: a | b
    assign y = ~(o & o);       // NOT of OR: four NANDs instead of one NOR
endmodule
\end{vcode}

\subsection{Example 3. De Morgan's Law}

Two forms of one function and a ``mismatch'' output that must always be zero.

\vfile{\codepath/b2_gates/demorgan.v}{code/b2\_gates/demorgan.v}

\textbf{How to read it.} The outputs \code{y1} and \code{y2} are the left and
right sides of the first law, \code{z1} and \code{z2} of the second. The
exclusive OR of two equal signals is 0, so \code{mismatch} $=0$ means both
pairs agree. The \code{demorgan} example from the Logra library uses the same
``mismatch'' output.

\textbf{Checking.} \code{tb\_demorgan.v} makes sure that \code{mismatch} is zero
on all four combinations. In Logra, sweep the inputs on the \emph{Waveforms}
tab: the \code{mismatch} line stays at zero. On the \emph{RTL} tab the
synthesizer may reduce both forms to one circuit, and that is fine: the
functions are equal.

\textbf{Finish it yourself.} Check the law for three variables:
$\overline{abc}=\overline a+\overline b+\overline c$.

\textbf{Sample answer.}\par\nopagebreak
\begin{vcode}
module demorgan3 (
    input  a, b, c,
    output mismatch      // 0 for all inputs
);
    wire y1 = ~(a & b & c);
    wire y2 = ~a | ~b | ~c;
    assign mismatch = y1 ^ y2;
endmodule
\end{vcode}

\subsection{Example 4. Exclusive OR in Three Forms}

Exclusive OR can be written with the operator \code{\^{}}, as ``different'', and
as ``at least one, but not both''. All three forms are one function.

\vfile{\codepath/b2_gates/xor_forms.v}{code/b2\_gates/xor\_forms.v}

\textbf{How to read it.}
\begin{itemize}
  \item \code{(a \& \textasciitilde b) | (\textasciitilde a \& b)} reads ``a without b, or b without a'': the standard form through AND, OR, NOT (Problem~5 of the chapter).
  \item \code{(a | b) \& \textasciitilde(a \& b)} is OR without the case ``both''.
  \item \code{same = \textasciitilde(a \^{} b)} is exclusive NOR: the equality of two bits (it will be needed for the comparator in Chapter~\ref{ch:arith}).
\end{itemize}
\textbf{Checking.} \code{tb\_xor\_forms.v} compares the three forms on all
combinations and separately checks the row $a=b=1$: exclusive OR gives 0 and OR
gives 1, the very difference that the ``Do it yourself'' task suggested finding.
In Logra look at \emph{RTL}: the synthesizer may reduce all three forms to one
circuit, because the functions are equal.

\textbf{Finish it yourself.} Make an equality comparator for two two-bit
numbers $a$ and $b$ using XNOR.

\textbf{Sample answer.}\par\nopagebreak
\begin{vcode}
module eq2 (
    input  [1:0] a, b,
    output       eq
);
    // the numbers are equal if both bits are equal
    assign eq = ~(a[0] ^ b[0]) & ~(a[1] ^ b[1]);
endmodule
\end{vcode}

\subsection{Example 5. The Security Panel Rule and Three-Input Gates}

The rule ``the door is open \emph{and} (it is night \emph{or} the arming key is
on)'' is written as one line; next to it are other functions of three
variables.

\vfile{\codepath/b2_gates/three_inputs.v}{code/b2\_gates/three\_inputs.v}

\textbf{How to read it.}
\begin{itemize}
  \item \code{alarm = door \& (night | armed)}: the parentheses are required. In Verilog, as in algebra, AND binds tighter than OR, and without them we would get $(\text{door}\cdot\text{night}) + \text{armed}$.
  \item \code{led\_n = \textasciitilde alarm}: the board LED lights at zero.
  \item \code{all3}, \code{odd3}, \code{two\_of\_3} are AND, exclusive OR (an odd number of ones) and ``at least two ones'' on three inputs; the last one will be useful for the carry of the adder in Chapter~\ref{ch:arith}.
\end{itemize}
\textbf{Checking.} \code{tb\_three\_inputs.v} sweeps the eight combinations and
checks that the alarm is 1 in exactly three rows (\code{101}, \code{110},
\code{111}). In Logra, sweep the inputs on the \emph{Waveforms} tab; the
\code{led\_n} LED goes out exactly where \code{alarm} is zero.

\textbf{Finish it yourself.} Add a panic button \code{panic}: when \code{panic}
$=1$ the alarm is always on.

\textbf{Sample answer.}\par\nopagebreak
\begin{vcode}
module alarm_panic (
    input  door, night, armed, panic,
    output alarm, led_n
);
    assign alarm = panic | (door & (night | armed));
    assign led_n = ~alarm;
endmodule
\end{vcode}

\section{Common Mistakes}

\begin{warnbox}[Common mistakes]
\begin{itemize}
  \item \textbf{Confusing OR and exclusive OR.} For $a=b=1$ OR gives 1, while
    exclusive OR gives 0. In everyday speech ``or'' often means exactly
    ``one of the two''.
  \item \textbf{Extra inversions.} $\overline{A+B}\neq\overline{A}+\overline{B}$:
    by De Morgan's law the inversion changes both the operation and the variables,
    $\overline{A+B}=\overline{A}\cdot\overline{B}$.
  \item \textbf{Thinking that nothing changes for a NAND with two connected inputs.}
    By connecting both inputs you get not a NAND but a NOT.
  \item \textbf{Operator precedence.} $A+B\cdot C$ means $A+(B\cdot C)$, not
    $(A+B)\cdot C$: AND binds more strongly than OR, just as multiplication binds more strongly than addition.
  \item \textbf{Verilog: forgetting the parentheses.} \code{door \& night |
    armed} is $(\text{door}\cdot\text{night})+\text{armed}$, not the rule you
    need. Always add parentheses when you mix \code{\&} and \code{|}.
  \item \textbf{Verilog: \code{\&\&} instead of \code{\&} and \code{!} instead
    of \code{\textasciitilde}.} For single bits the result is the same, but on
    buses \code{\&\&} and \code{!} give a single ``not zero'' bit, while
    \code{\&} and \code{\textasciitilde} work on every bit separately.
  \item \textbf{Verilog: two \code{assign} statements for one wire.} If the same
    \code{wire} receives a value from two expressions, it has two sources; the
    synthesizer reports an error (or the result is the undefined value
    \code{x}).
\end{itemize}
\end{warnbox}

\section{Questions and Problems}
\label{sec:tasks-gates}

The answers are collected in Appendix~\ref{app:answers}.

\begin{taskbox}[Chapter problems]
\begin{enumerate}
  \item \stars{1}~Build the truth table of the function $F = A\cdot B + \overline{A}\cdot C$.
  \item \stars{1}~Simplify the expression $A + A\cdot B$ and name the law you used.
  \item \stars{2}~Apply De Morgan's law and write $\overline{A + B\cdot C}$ without the common inversion over the sum.
  \item \stars{2}~Build the NOT, AND and OR gates using only NAND gates. How many gates are needed for each?
  \item \stars{2}~Express exclusive OR through AND, OR and NOT.
  \item \stars{1}~How many rows does the truth table of a function of six variables have?
  \item \stars{2}~\emph{Find the mistake.} A student ``simplified'' $\overline{A+B}=\overline{A}+\overline{B}$. Find an input combination for which the equality fails, and write the correct law.
  \item \stars{2}~\emph{Find the mistake.} A student replaced an exclusive OR gate in a circuit with an OR gate because ``they do the same thing''. For which input combination will the behavior of the circuit change?
  \item \stars{3}~\emph{Project.} Stage 2 of the ``Security Panel'' project: $\text{alarm}=\text{door}\cdot(\text{night}+\text{armed})$. Build the truth table, build the circuit in Logra and send the output to the board LED (it lights on zero, so an inversion is needed). Acceptance criterion: on the \emph{Waveforms} tab, when going through all eight combinations, the LED is on in exactly three.
  \item \stars{3}~\emph{Project.} Build the panel output \code{led\_n} using only NAND gates. How many gates are needed? Acceptance criterion: the waveform matches the waveform of the circuit from problem~9 on all combinations.
\end{enumerate}
\end{taskbox}

\section{What's Next}

We can write logic functions and build them from gates, on paper and in the
Logra Designer. But a circuit of ten gates can still be drawn with the mouse,
while one of a hundred cannot. How do we write a circuit as \emph{text}, so that
Logra turns it into gates by itself? The answer is the Verilog language: in
Chapter~\ref{ch:verilog} we describe the same gates in it and immediately run
the result on the RTL and Waveforms tabs. After that, in Chapter~\ref{ch:comb},
we return to the security panel with three sensors: from a verbal condition we
go to a truth table and a circuit, and study minimization along the way.
